%%
%% part3-customization/ch20-layout.tex
%%
%% Chapter 20: Page Layout — Geometry, margins, line
%% spacing, headers, footers, and paragraph settings.
%%

\chapter{صفحه‌آرایی}
\label{chap:layout}

\section{چرا این موضوع مهم است؟}
\label{sec:layout-why}

صفحه‌آرایی، پایه‌ی کیفیت بصری یک کتاب است.
حاشیه‌های درست، فاصله‌ی خطوط مناسب، سرصفحه‌ی
خوانا و پاصفحه‌ی تمیز، همگی به خوانایی و
زیبایی کتاب کمک می‌کنند.

\lr{MatinBook} یک چیدمان حرفه‌ای بر اساس
استاندارد \lr{Boyer/Stewart} ارائه می‌دهد.
در این فصل، با اجزای این چیدمان آشنا
می‌شوید و یاد می‌گیرید که چگونه آن را
شخصی‌سازی کنید.

\mnote{چیدمان \lr{MatinBook} از کتاب‌های درسی دانشگاهی مانند \lr{Boyer} و \lr{Stewart} الهام گرفته شده است.}

\section{هندسه‌ی صفحه}
\label{sec:layout-geometry}

\subsection{ابعاد کاغذ}
\label{sec:layout-paper}

\lr{MatinBook} از ابعاد \pnum{۱۷.۸} ×
\pnum{۲۵.۴} سانتی‌متر (نزدیک به \lr{7 × 10}
اینچ) استفاده می‌کند. این ابعاد، استاندارد
نشر آکادمیک بین‌المللی است.

\begin{latin}
\begin{minted}{latex}
% |\pc{در فایل}|\lr{mb-layout.sty}:
\geometry{
    paperwidth=17.8cm,
    paperheight=25.4cm,
    ...
}
\end{minted}
\end{latin}

\mnote{این ابعاد، نزدیک به نسبت طلایی است و برای کتاب‌های فنی، خوانایی بالایی دارد.}

\subsection{حاشیه‌ها در \lr{mainmatter}}
\label{sec:layout-margins-main}

در متن اصلی، \lr{MatinBook} از حاشیه‌ی
بیرونی پهن (\pnum{۵.۹} سانتی‌متر) استفاده
می‌کند:

\begin{latin}
\begin{minted}{latex}
\geometry{
    paperwidth=17.8cm,
    paperheight=25.4cm,
    top=1.7cm,
    bottom=1.7cm,
    inner=1.5cm,
    outer=5.9cm,              % |\pc{حاشیه‌ی پهن}|%
    marginparwidth=3.5cm,     % |\pc{عرض یادداشت}|%
    marginparsep=0.6cm,       % |\pc{فاصله از متن}|%
    headheight=15pt,
    headsep=5pt,
    footskip=18pt,
    bindingoffset=0.4cm
}
\end{minted}
\end{latin}

جدول \cref{tab:layout-margins} این مقادیر
را خلاصه می‌کند.

\begin{matintable}{حاشیه‌های متن اصلی}{tab:layout-margins}
\begin{tblr}{
    width = \textwidth,
    colspec = {l X[c] X[l]},
    hlines, vlines,
    colsep = 8pt,
    rowsep = 4pt,
    row{1} = {font=\bfseries},
}
حاشیه & مقدار & کاربرد \\
بالا & \pnum{۱.۷} سانتی‌متر & فاصله از لبه‌ی بالا \\
پایین & \pnum{۱.۷} سانتی‌متر & فاصله از لبه‌ی پایین \\
داخلی & \pnum{۱.۵} سانتی‌متر & نزدیک عطف \\
بیرونی & \pnum{۵.۹} سانتی‌متر & برای یادداشت حاشیه \\
عرض یادداشت & \pnum{۳.۵} سانتی‌متر & ناحیه‌ی یادداشت \\
فاصله از متن & \pnum{۰.۶} سانتی‌متر & بین متن و یادداشت \\
فضای صحافی & \pnum{۰.۴} سانتی‌متر & برای صحافی \\
\end{tblr}
\end{matintable}

\mnote{عرض متن در این حالت، حدود \pnum{۱۰.۳} سانتی‌متر است که برای متن فارسی با فونت \lr{XB Niloofar} بهینه است.}

\subsection{حاشیه‌ها در \lr{frontmatter} و \lr{backmatter}}
\label{sec:layout-margins-front}

در پیش‌متن و پس‌متن، هندسه‌ی صفحه متقارن
می‌شود (بدون حاشیه‌ی پهن):

\begin{latin}
\begin{minted}{latex}
% |\pc{در فایل}|\lr{mb-layout.sty}:
\newcommand{\frontmattergeometry}{%
    \newgeometry{
        paperwidth=17.8cm,
        paperheight=25.4cm,
        top=1.7cm,
        bottom=1.7cm,
        inner=1.5cm,
        outer=1.5cm,          % |\pc{حاشیه‌ی متقارن}|%
        marginparwidth=0pt,
        ...
    }%
}
\end{minted}
\end{latin}

\mnote{دلیل حذف حاشیه‌ی پهن در پیش‌متن و پس‌متن این است که این بخش‌ها به یادداشت حاشیه‌ای نیاز ندارند.}

\section{فاصله‌ی خطوط}
\label{sec:layout-spacing}

\lr{MatinBook} از فاصله‌ی خطوط \pnum{۱.۱۵}
استفاده می‌کند که مطابق سبک کتاب‌های درسی
\lr{Boyer/Stewart} است:

\begin{latin}
\begin{minted}{latex}
\setstretch{1.15}
\end{minted}
\end{latin}

\mnote{فاصله‌ی \pnum{۱.۱۵} برای متن فارسی با فونت \lr{XB Niloofar} بهینه است و خوانایی را افزایش می‌دهد.}

\subsection{تنظیمات پاراگراف}
\label{sec:layout-paragraph}

\begin{latin}
\begin{minted}{latex}
\setlength{\parindent}{1em}
\setlength{\parskip}{0pt}
\end{minted}
\end{latin}

\begin{itemize}
    \item \cmd{parindent}: تورفتگی پاراگراف.
    \item \cmd{parskip}: فاصله‌ی بین پاراگراف‌ها.
\end{itemize}

\mnote{در فارسی، تورفتگی \lr{۱em} استاندارد است. پاراگراف اول هر بخش، تورفتگی ندارد.}

\section{سرصفحه و پاصفحه}
\label{sec:layout-headers}

\lr{MatinBook} از \lr{fancyhdr} برای مدیریت
سرصفحه و پاصفحه استفاده می‌کند.

\subsection{صفحات فرد و زوج}
\label{sec:layout-headers-oddeven}

\begin{itemize}
    \item \textbf{صفحات فرد:} نام فصل در بالا
    راست، شماره صفحه در پایین راست.
    \item \textbf{صفحات زوج:} نام فصل در بالا
    چپ، شماره صفحه در پایین چپ.
\end{itemize}

\begin{latin}
\begin{minted}{latex}
\pagestyle{fancy}
\fancyhf{}

\fancyhead[LE]{%
    \ifnum\c@chapter>0
        \footnotesize\textbf{\leftmark}%
    \fi
}

\fancyhead[RO]{%
    \ifnum\c@chapter>0
        \footnotesize\textbf{\leftmark}%
    \fi
}

\fancyfoot[LE,RO]{\footnotesize\textbf{\thepage}}
\end{minted}
\end{latin}

\mnote{در \lr{MatinBook}، از \lr{\textbackslash leftmark} برای هر دو سرصفحه استفاده می‌شود. یعنی نام فصل در هر دو صفحه نمایش داده می‌شود.}

\subsection{صفحات شروع فصل}
\label{sec:layout-headers-plain}

صفحات شروع فصل، از استایل \lr{plain} استفاده
می‌کنند و فقط شماره صفحه را در پاصفحه
نمایش می‌دهند:

\begin{latin}
\begin{minted}{latex}
\fancypagestyle{plain}{%
    \fancyhf{}
    \fancyfoot[LE,RO]{\footnotesize\textbf{\thepage}}
    \renewcommand{\headrulewidth}{0pt}
    \renewcommand{\footrulewidth}{0pt}
}
\end{minted}
\end{latin}

\mnote{این طراحی تمیز، توجه خواننده را به ابتدای فصل جلب می‌کند و فضای بیشتری برای عنوان فصل فراهم می‌آورد.}

\subsection{خط جداکننده‌ی سرصفحه}
\label{sec:layout-headers-rule}

\begin{latin}
\begin{minted}{latex}
\renewcommand{\headrulewidth}{0.4pt}
\renewcommand{\footrulewidth}{0pt}
\end{minted}
\end{latin}

\mnote{خط جداکننده‌ی سرصفحه با ضخامت \pnum{۰.۴pt} طراحی شده است. پاصفحه خط ندارد.}

\section{تنظیمات کپشن}
\label{sec:layout-captions}

\begin{latin}
\begin{minted}{latex}
\captionsetup{
    font=footnotesize,
    labelfont=bf,
    skip=8pt,
    justification=centering
}
\end{minted}
\end{latin}

\mnote{کپشن‌ها با فونت \lr{footnotesize} و برچسب پررنگ نمایش داده می‌شوند.}

\section{تنظیمات پانویس}
\label{sec:layout-footnotes}

\begin{latin}
\begin{minted}{latex}
\setlength{\footnotesep}{8pt}
\renewcommand{\footnoterule}{%
    \kern-3pt
    \hrule width 0.3\textwidth height 0.4pt
    \kern2.6pt
}
\end{minted}
\end{latin}

\mnote{خط جداکننده‌ی پانویس، ۳۰٪ عرض متن است. این طراحی، استاندارد کتاب‌های آکادمیک است.}

\section{حالت \lr{draft}}
\label{sec:layout-draft}

در حالت \lr{draft}، \lr{MatinBook} واترمارک
\lr{DRAFT} را فعال می‌کند:

\begin{latin}
\begin{minted}{latex}
\if@matin@draft
    \AtBeginDocument{%
        \RequirePackage{draftwatermark}%
        \SetWatermarkText{\lr{DRAFT}}%
        \SetWatermarkScale{1.5}%
        \SetWatermarkColor[gray]{0.9}%
    }
\fi
\end{minted}
\end{latin}

\mnote{واترمارک در حالت \lr{draft} به‌عنوان یادآوری برای بازگشت به حالت \lr{final} عمل می‌کند.}

\section{شخصی‌سازی چیدمان}
\label{sec:layout-customization}

برای شخصی‌سازی چیدمان، فایل
\file{mb-layout.sty} را ویرایش کنید. سه
راه اصلی وجود دارد:

\begin{enumerate}
    \item \textbf{تغییر هندسه:} مقادیر
    \cmd{geometry} را تغییر دهید.
    \item \textbf{تغییر سرصفحه:} تنظیمات
    \cmd{fancyhdr} را تغییر دهید.
    \item \textbf{تغییر فاصله‌ها:} مقادیر
    \cmd{setstretch} و \cmd{parindent} را
    تغییر دهید.
\end{enumerate}

\mnote{تغییر هندسه‌ی صفحه، بر کل کتاب تأثیر می‌گذارد. قبل از تغییر، از کتاب پشتیبان بگیرید.}

\section{خطاهای رایج}
\label{sec:layout-mistakes}

\subsection{فراموش کردن \cmd{frontmattergeometry}}
\label{sec:layout-mistake-frontgeo}

\textbf{مشکل:} اگر \cmd{frontmattergeometry}
را فراخوانی نکنید، پیش‌متن با حاشیه‌ی پهن
نمایش داده می‌شود.

\textbf{راه‌حل:} همیشه پس از \cmd{frontmatter}،
دستور \cmd{frontmattergeometry} را فراخوانی
کنید.

\subsection{فراموش کردن \cmd{mainmattergeometry}}
\label{sec:layout-mistake-maingeo}

\textbf{مشکل:} اگر \cmd{mainmattergeometry}
را فراخوانی نکنید، متن اصلی با حاشیه‌ی
متقارن نمایش داده می‌شود.

\textbf{راه‌حل:} همیشه پس از \cmd{mainmatter}،
دستور \cmd{mainmattergeometry} را فراخوانی
کنید.

\subsection{سرریز متن به حاشیه}
\label{sec:layout-mistake-overflow}

\textbf{مشکل:} اگر متن یا شکل از حاشیه سرریز
کند، خطای \lr{Overfull hbox} می‌دهد.

\textbf{راه‌حل:} از \cmd{emergencystretch} و
\cmd{tolerance} استفاده کنید (که در
\lr{MatinBook} از پیش تنظیم شده است).

\section{تمرین‌ها}
\label{sec:layout-exercises}

\begin{exercise}
    حاشیه‌ی بیرونی را از \pnum{۵.۹} به
    \pnum{۶.۵} سانتی‌متر تغییر دهید و کتاب
    را دوباره کامپایل کنید.
    \label{exc:layout-margins}
\end{exercise}

\begin{exercise}
    فاصله‌ی خطوط را از \pnum{۱.۱۵} به
    \pnum{۱.۳} تغییر دهید و نتیجه را ببینید.
    \label{exc:layout-spacing}
\end{exercise}

\begin{exercise}
    رنگ سرصفحه را تغییر دهید.
    \label{exc:layout-headers}
\end{exercise}

\begin{exercise}
    کتاب را در حالت \lr{draft} کامپایل کنید
    و واترمارک را ببینید.
    \label{exc:layout-draft}
\end{exercise}

\section{خلاصه}
\label{sec:layout-summary}

در این فصل، با صفحه‌آرایی \lr{MatinBook}
آشنا شدیم:

\begin{itemize}
    \item ابعاد کاغذ \pnum{۱۷.۸} ×
    \pnum{۲۵.۴} سانتی‌متر.

    \item حاشیه‌ی بیرونی پهن \pnum{۵.۹}
    سانتی‌متر در متن اصلی.

    \item حاشیه‌ی متقارن \pnum{۱.۵} سانتی‌متر
    در پیش‌متن و پس‌متن.

    \item فاصله‌ی خطوط \pnum{۱.۱۵}.

    \item تورفتگی پاراگراف \pnum{۱em}.

    \item سرصفحه با نام فصل، پاصفحه با شماره
    صفحه.

    \item خط جداکننده‌ی سرصفحه با ضخامت
    \pnum{۰.۴pt}.

    \item تنظیمات کپشن و پانویس.

    \item حالت \lr{draft} با واترمارک.

    \item خطاهای رایج شامل فراموش کردن
    \cmd{frontmattergeometry}، فراموش کردن
    \cmd{mainmattergeometry}، و سرریز متن
    است.
\end{itemize}

در فصل بعدی (\cref{chap:themes})، با تم‌ها
در \lr{MatinBook} آشنا می‌شوید و یاد
می‌گیرید که چگونه تم سفارشی بسازید.